\section{Đạo đức và dữ liệu}\label{sec:ethics}

\paragraph{Dữ liệu.} Không có dữ liệu của trẻ thật trong các thực nghiệm; học sinh là kịch bản và được ẩn danh (S1). Các tài liệu nguồn do giáo viên biên soạn có hỗ trợ của AI, không sao chép từ sách của nhà xuất bản, và không được công bố; bài báo công bố thống kê tổng hợp và các số liệu sinh bảng biểu, không công bố nguyên văn tài liệu. Khi hệ thống chạy với học sinh thật, chỉ tên gọi (không họ) được gửi cho dịch vụ mô hình, câu kiểm tra không bao giờ được gửi, và các khâu pháp lý, đồng ý của phụ huynh thuộc trách nhiệm của đơn vị triển khai.

\paragraph{An toàn.} Gia sư không hỏi thông tin cá nhân (bị chặn bởi bộ kiểm chứng), không khen trí thông minh, và với lời đáng lo (bị bắt nạt, muốn làm đau mình) khuyên nói với thầy cô hoặc bố mẹ và đánh dấu để giáo viên thấy. Lời học sinh được đối xử như dữ liệu, không phải chỉ thị, nên lời ``bỏ qua mọi hướng dẫn'' không đổi hành vi. Các cơ chế này chưa được kiểm trong thực địa.

\paragraph{Vai trò của AI.} Trợ lý AI được dùng để viết mã, soạn tập câu hỏi và nhãn của chúng, duyệt nội dung theo ủy quyền của giáo viên, và soạn bản thảo của bài báo này; mọi con số trong bài được sinh bằng mã từ các tệp dữ liệu công bố kèm, và chủ dự án (đồng thời là giáo viên kiểm soát) duyệt mọi đặc tả sư phạm. Chúng tôi nêu điều này vì nó là một phần của mối đe dọa đến tính hợp lệ ở Mục~\ref{sec:limits}.

\paragraph{Sử dụng.} CAVAT là công cụ hỗ trợ giáo viên, không thay thế giáo viên. Các con số của bài không đủ để biện minh cho việc triển khai không giám sát với trẻ em; các thử nghiệm với học sinh thật cần có thiết kế, đồng ý và giám sát riêng.
